
\begin{obeactivity}
\textbf{Aktivitas 4.1: Praktik Analisis Frekuensi}
1. Gunakan alat \textit{online} seperti \url{https://www.cryptool.org/en/cto/n-gram-analysis}.
2. Masukkan sebuah cipherteks yang dihasilkan dari cipher abjad-tunggal.
3. Analisis distribusi frekuensi hurufnya, lalu petakan huruf yang paling sering muncul ke huruf \code{E} atau \code{A}.
4. Lakukan iterasi \textit{trial-and-error} sampai pesan dapat terbaca.
\end{obeactivity}

Langkah-langkah pada Aktivitas 4.1 dapat dijalankan melalui program pada
Listing~\ref{lst:caesar-frekuensi}. Program itu mengenkripsi pesan, menghitung
frekuensi huruf cipherteks, lalu menebak kunci dengan menyelaraskan huruf
tersering ke urutan huruf bahasa Indonesia.

\lstinputlisting[
    language=Python,
    caption={Cipher Caesar: enkripsi, dekripsi, dan pemecahan dengan analisis frekuensi},
    label={lst:caesar-frekuensi}
]{\codefile{python/caesar.py}}

Program yang sama tersedia dalam C++ (\texttt{code/cpp/caesar.cpp}) dan MATLAB
(\texttt{code/matlab/caesar.m}).

\begin{obeactivity}
\textbf{Aktivitas 4.2: Kriptanalisis Vigen\`ere dengan Metode Kasiski}
1. Bangkitkan cipherteks Vigen\`ere dari kalimat bahasa Indonesia yang panjang (minimal 200 huruf) dengan kunci berpanjang 4 sampai 6 huruf. Gunakan alat \textit{online} seperti \url{https://www.dcode.fr/vigenere-cipher} atau \url{https://cryptii.com/pipes/vigenere-cipher}.
2. Tanpa melihat kunci, cari kriptogram yang berulang dan catat jarak antar kemunculannya.
3. Faktorkan setiap jarak, lalu tentukan panjang kunci $m$ sebagai faktor yang muncul pada semua himpunan faktor.
4. Bagi cipherteks menjadi $m$ kelompok menurut posisi modulo $m$, lalu jalankan analisis frekuensi pada setiap kelompok untuk menebak huruf kunci.
5. Bandingkan kunci yang diperoleh dengan kunci aslinya.
\end{obeactivity}

\begin{obeactivity}
\textbf{Aktivitas 4.3: Membandingkan Cipher melalui Indeks Kebetulan}
1. Enkripsikan satu kalimat yang sama dengan Caesar cipher dan Vigen\`ere cipher.
2. Hitung frekuensi huruf setiap cipherteks.
3. Hitung indeks kebetulan dengan rumus pada Bagian~\ref{sec:indeks-kebetulan}.
4. Bandingkan hasilnya dengan nilai acuan $0{,}065$ untuk cipher abjad-tunggal dan $0{,}038$ untuk teks acak.
5. Simpulkan cipher mana yang distribusinya lebih mendekati teks acak.
\end{obeactivity}
